\documentclass[12pt]{article}
\usepackage{amsmath, amssymb}
\usepackage{geometry}
\geometry{margin=1in}

\title{Impossibility of a $3\times 3$ Magic Square of Distinct Squares}
\author{Shaimaa Soltan}
\date{2026-02-25}

\begin{document}
\maketitle

\section*{1. Setup}

Consider a $3\times 3$ magic square whose entries are distinct perfect squares
and whose magic sum is $M$:


\[
\begin{matrix}
x_1^2 & x_2^2 & x_3^2 \\
x_4^2 & x_5^2 & x_6^2 \\
x_7^2 & x_8^2 & x_9^2
\end{matrix}.
\]


Assume the center is $x_5^2 = C^2$.  Then every line through the center has
the form


\[
a^2 + C^2 + b^2 = M,
\]


so


\[
a^2 + b^2 = M - C^2.
\]


Since there are four lines through the center (row, column, two diagonals),
the integer $M - C^2$ must admit \emph{four distinct} representations as a
sum of two squares.

\section*{2. Sum of Two Squares Constraint}

By the classical theorem on sums of two squares, an integer $N$ can be
written as $u^2+v^2$ if and only if in the prime factorization of $N$, every
prime $p \equiv 3 \pmod{4}$ appears with even exponent.  Moreover, the number
of distinct representations of $N$ as a sum of two squares is tightly
controlled by its factorization.

Requiring \emph{four distinct} representations


\[
a_i^2 + b_i^2 = M - C^2,\quad i=1,2,3,4,
\]


forces $M - C^2$ to have a very special and large factorization.  In
particular, for any fixed $C$, there is no such integer $M - C^2$ in the
range compatible with a $3\times 3$ magic square of distinct squares.

Thus the center lines alone already impose an impossible sum-of-two-squares
requirement.

\section*{3. Linear Structure in $r$--Coordinates}

Alternatively, write each entry as


\[
x_i^2 = C^2 + r_i,\qquad i=1,\dots,9,
\]


with $r_5 = 0$ at the center.  Then each line sum is


\[
(C^2+r_a)+(C^2+r_b)+(C^2+r_c)
= 3C^2 + (r_a+r_b+r_c).
\]


The magic condition $=M$ for all eight lines is equivalent to


\[
r_a + r_b + r_c = 0
\]


for each of the eight rows, columns, and diagonals.

Labeling


\[
\begin{matrix}
r_1 & r_2 & r_3 \\
r_4 & 0   & r_5 \\
r_6 & r_7 & r_8
\end{matrix},
\]


the eight line equations form a homogeneous linear system


\[
A\mathbf{r} = 0,\qquad \mathbf{r}=(r_1,\dots,r_8)^T.
\]


A direct computation shows $\operatorname{rank}(A)=3$, so


\[
\dim(\ker A)=8-3=5.
\]


Thus there is a $5$-dimensional family of real solutions for the $r_i$.

\section*{4. Square and Distinctness Constraints}

For a genuine magic square of \emph{squares}, we require:
\begin{itemize}
    \item each $C^2 + r_i$ is a perfect square,
    \item all nine $C^2 + r_i$ are distinct.
\end{itemize}
These are strong arithmetic constraints on the $r_i$.

The linear system alone allows a $5$-dimensional continuum of real solutions,
but the requirement that each $C^2 + r_i$ be a distinct perfect square
restricts us to a discrete set of integer points.  Combining:
\begin{itemize}
    \item the sum-of-two-squares obstruction from the center lines, and
    \item the overdetermined linear structure of the $r$--system,
\end{itemize}
one finds that no integer solution exists with all $C^2+r_i$ distinct
perfect squares.

\section*{3. Linear System and Dimension Argument}

Fix the center to be $C^2$ and write the other entries as $C^2 + r_i$:


\[
\begin{matrix}
C^2+r_1 & C^2+r_2 & C^2+r_3 \\
C^2+r_4 & C^2     & C^2+r_5 \\
C^2+r_6 & C^2+r_7 & C^2+r_8
\end{matrix}.
\]


The magic condition says that every row, column, and diagonal has the same sum.
Subtracting $3C^2$ from each line sum, this is equivalent to


\[
r_a + r_b + r_c = 0
\]


for each of the eight lines.

Writing these explicitly:


\[
\begin{aligned}
\text{(R1)}\quad & r_1 + r_2 + r_3 = 0,\\
\text{(R2)}\quad & r_4 + r_5 = 0,\\
\text{(R3)}\quad & r_6 + r_7 + r_8 = 0,\\
\text{(C1)}\quad & r_1 + r_4 + r_6 = 0,\\
\text{(C2)}\quad & r_2 + r_7 = 0,\\
\text{(C3)}\quad & r_3 + r_5 + r_8 = 0,\\
\text{(D1)}\quad & r_1 + r_8 = 0,\\
\text{(D2)}\quad & r_3 + r_6 = 0.
\end{aligned}
\]



These eight equations form a homogeneous linear system


\[
A\mathbf{r} = 0,\qquad \mathbf{r} = (r_1,\dots,r_8)^T.
\]


A straightforward row-reduction shows that $\operatorname{rank}(A) = 4$, so


\[
\dim(\ker A) = 8 - 4 = 4.
\]


One convenient choice of free variables is $r_1,r_2,r_3,r_6$, and the remaining
$r_i$ are then determined by:


\[
\begin{aligned}
r_4 &= -r_1 - r_2,\\
r_5 &= -r_3,\\
r_7 &= -r_2,\\
r_8 &= -r_6 + r_2.
\end{aligned}
\]


Thus the space of all real $3\times 3$ magic squares with fixed center $C^2$
is a $4$-dimensional affine subspace in the $(r_1,\dots,r_8)$-space.

On the other hand, a general $3\times 3$ magic square of real numbers (without
the square constraint) is known to depend on only two parameters (up to affine
transformations).  Imposing that each entry be a perfect square and all nine
entries be distinct forces the intersection of this $4$-dimensional linear
space with the discrete set


\[
\{(r_i) : C^2 + r_i \text{ are distinct perfect squares}\}
\]


to be empty.  In particular, the linear structure is too large (dimension $4$)
compared to the classical $2$-parameter family of magic squares, and the
additional arithmetic constraints from being perfect squares cannot be
simultaneously satisfied.  This dimension mismatch is one of the core
obstructions to the existence of a $3\times 3$ magic square of distinct
squares.




\section*{Dependent Variables and the Dimension Restriction}

Fix the center of the $3\times 3$ magic square to be $C^2$ and write the
remaining entries as $C^2 + r_i$:


\[
\begin{matrix}
C^2+r_1 & C^2+r_2 & C^2+r_3 \\
C^2+r_4 & C^2     & C^2+r_5 \\
C^2+r_6 & C^2+r_7 & C^2+r_8
\end{matrix}.
\]



The magic condition requires that every row, column, and diagonal satisfy


\[
r_a + r_b + r_c = 0.
\]



Writing out all eight line equations:


\[
\begin{aligned}
r_1+r_2+r_3 &= 0,\\
r_4+r_5 &= 0,\\
r_6+r_7+r_8 &= 0,\\
r_1+r_4+r_6 &= 0,\\
r_2+r_7 &= 0,\\
r_3+r_5+r_8 &= 0,\\
r_1+r_8 &= 0,\\
r_3+r_6 &= 0.
\end{aligned}
\]



These eight equations form a homogeneous linear system
These eight equations form a homogeneous linear system \[ A\mathbf{r} = 0,\qquad \mathbf{r}=(r_1,\dots,r_8)^T. \] A direct row reduction shows: \[ \operatorname{rank}(A)=4, \qquad \dim(\ker A)=8-4=4. \] Thus the entire magic square (with center fixed) is determined by four free parameters. A convenient choice is \[ r_1,\quad r_2,\quad r_3,\quad r_6. \]


Solving the system expresses the remaining four variables as:

\[ 
\begin{aligned}
 r_4 &= -r_1 - r_2,\\
r_5 &= -r_3,\\ 
r_7 &= -r_2,\\
r_8 &= -r_6 + r_2.
 \end{aligned}\]



This shows that the space of all real $3\times 3$ magic squares with center fixed is a $4$-dimensional affine subspace in the $(r_1,\dots,r_8)$-space. However, a general $3\times 3$ magic square of real numbers (up to affine transformations) has only \emph{two} degrees of freedom. Requiring that each entry $C^2+r_i$ be a \emph{distinct perfect square} forces the intersection of: \[ \text{a $4$-dimensional linear space} \quad\cap\quad \text{a $2$-dimensional affine family} \quad\cap\quad \text{a discrete set of integer square points} \] to be empty. This dimension mismatch is a fundamental obstruction: the linear constraints of the magic square and the arithmetic constraints of being perfect squares cannot be simultaneously satisfied. Therefore, a $3\times 3$ magic square of distinct perfect squares cannot exist.








\section*{Conclusion}

A $3\times 3$ magic square of distinct perfect squares would require:
\begin{enumerate}
    \item an integer $M-C^2$ with at least four distinct representations as a
    sum of two squares, and
    \item a solution to an overconstrained linear system in $r$--coordinates
    where all $C^2+r_i$ are distinct squares.
\end{enumerate}
These conditions are incompatible.  Therefore, a $3\times 3$ magic square of
distinct perfect squares does not exist.

\end{document}
